\section{Proof Architecture}
\label{sec:proof}

Our formalization is organized into five independent tracks that compose to yield
the main theorem and its extensions. \Cref{fig:dag} shows the dependency
structure. We present the mathematical arguments here; the Lean formalization
is discussed in \cref{sec:formalization}.

\begin{figure}[ht]
\centering
\small
\[
\underbrace{
  \begin{array}{c}
  \texttt{Telescoping} \\[2pt]
  X^n - Y^n = \sum X^k(X{-}Y)Y^{n{-}1{-}k}
  \end{array}
}_{\text{Track 1}}
\quad
\underbrace{
  \begin{array}{c}
  \texttt{ExpBounds} \\[2pt]
  \|e^a\| \le e^{\|a\|}
  \end{array}
}_{\text{Track 2}}
\quad
\underbrace{
  \begin{array}{c}
  \texttt{ExpDivPow} \\[2pt]
  (e^{a/n})^n = e^a
  \end{array}
}_{\text{Track 4}}
\]
\[
\downarrow \qquad\qquad\qquad\qquad \downarrow
\]
\[
\underbrace{
  \begin{array}{c}
  \texttt{StepError} \\[2pt]
  \|e^a e^b - e^{a+b}\| \le 2\|a\|\|b\|\,e^{\|a\|+\|b\|}
  \end{array}
}_{\text{Track 3 (depends on Track 2)}}
\]
\[
\downarrow
\]
\[
\boxed{
  \begin{array}{c}
  \texttt{Assembly} \\[2pt]
  (e^{A/n}\,e^{B/n})^n \to e^{A+B} \quad [O(1/n)]
  \end{array}
}
\]
\[
\swarrow \qquad\qquad\qquad \searrow
\]
\[
\underbrace{
  \begin{array}{c}
  \texttt{StrangSplitting} \\[2pt]
  O(1/n^2)
  \end{array}
}_{\text{Extensions}}
\quad
\underbrace{
  \begin{array}{c}
  \texttt{Multi\{Operator,Strang\}} \\[2pt]
  m\text{-operator}
  \end{array}
}_{\text{Extensions}}
\quad
\underbrace{
  \begin{array}{c}
  \texttt{CommutatorScaling} \\[2pt]
  \|[B,A]\|\text{-scaling}
  \end{array}
}_{\text{Track 5 (independent)}}
\]
\caption{Dependency structure of the formalization.
  Tracks~1--4 compose into the main theorem; extensions branch from the
  assembly or develop independently.}
\label{fig:dag}
\end{figure}

\subsection{Modular decomposition}
\label{sec:modular}

The proof of the Lie--Trotter formula decomposes into four independent
components:
\begin{description}[leftmargin=1.5em]
  \item[Track~1 (Algebraic):] The telescoping identity and its norm bound.
  \item[Track~2 (Analytic):] Exponential series remainder estimates.
  \item[Track~3 (Step error):] The single-step approximation bound
    $\norm{e^a\,e^b - e^{a+b}} \le 2\norm{a}\norm{b}\,e^{\norm{a}+\norm{b}}$.
  \item[Track~4 (Connecting):] The identity $(e^{a/n})^n = e^a$.
\end{description}
Tracks 1--4 are combined in an assembly step to yield the convergence rate,
from which the limit theorem follows by a standard $\varepsilon$-$N$ argument.
The commutator-scaling bounds (Track~5) are an independent extension that
refines the step error.

\subsection{Telescoping and norm bounds}
\label{sec:telescoping}

The algebraic identity
\begin{equation}\label{eq:telescoping}
  X^n - Y^n = \sum_{k=0}^{n-1} X^k\,(X-Y)\,Y^{n-1-k}
\end{equation}
is proved by induction on $n$. The key step is factoring $Y$ out of the inner
sum to match the inductive hypothesis: writing $Y^{m-k} = Y^{m-1-k} \cdot Y$
converts $\sum_{k<m} X^k(X-Y)Y^{m-k}$ into $(\sum_{k<m} X^k(X-Y)Y^{m-1-k}) \cdot Y$.

The norm bound follows by the triangle inequality applied to each summand:
\begin{equation}\label{eq:norm-telescoping}
  \norm{X^n - Y^n} \le n\,\norm{X-Y}\,M^{n-1},
  \qquad M = \max\{\norm{X}, \norm{Y}\}.
\end{equation}

\subsection{Exponential series estimates}
\label{sec:exp-bounds}

We establish four bounds on the exponential in a complete normed algebra:
\begin{align}
  \norm{e^a} &\le e^{\norm{a}},
    \label{eq:norm-exp} \\
  \norm{e^a - 1} &\le e^{\norm{a}} - 1,
    \label{eq:norm-exp-sub-one} \\
  e^x - 1 - x &\le \tfrac{x^2}{2}\,e^x
    \quad (x \ge 0),
    \label{eq:exp-quadratic} \\
  \norm{e^a - 1 - a} &\le \tfrac{\norm{a}^2}{2}\,e^{\norm{a}}.
    \label{eq:norm-exp-quadratic}
\end{align}

The proofs use the power series representation $e^a = \sum_{n=0}^{\infty} a^n/n!$.
Bounds \eqref{eq:norm-exp}--\eqref{eq:norm-exp-sub-one} follow from
$\norm{\sum a^n/n!} \le \sum \norm{a}^n/n! = e^{\norm{a}}$ with appropriate
index shifts. Bound \eqref{eq:exp-quadratic} uses the termwise estimate
$x^{n+2}/(n+2)! \le (x^2/2) \cdot x^n/n!$ via the factorial inequality
$2 \cdot n! \le (n+2)!$. Bound \eqref{eq:norm-exp-quadratic} combines
\eqref{eq:exp-quadratic} with the same tsum comparison technique as
\eqref{eq:norm-exp}.

\begin{remark}
  Bound \eqref{eq:norm-exp} was not available in \Mathlib for general
  Banach algebras; \Mathlib only had the specialized result
  \texttt{Complex.norm\_exp\_le\_exp\_norm} for $\mathbb{C}$.
  Our proof from the power series fills this gap.
\end{remark}

\subsection{Step error via algebraic factorization}
\label{sec:step-error}

The single-step bound is the most delicate component. Rather than expanding
$e^a\,e^b$ and $e^{a+b}$ to second order and bounding cross-terms directly
(which is tractable on paper but unwieldy in a proof assistant due to
non-commutative bookkeeping), we use an algebraic factorization:
\begin{equation}\label{eq:factorization}
  e^a\,e^b - e^{a+b}
    = (e^a - 1)(e^b - 1) - \bigl(e^{a+b} - e^a - e^b + 1\bigr).
\end{equation}
The first term is bounded by $(e^{\norm{a}}-1)(e^{\norm{b}}-1)$ via
\eqref{eq:norm-exp-sub-one}. For the second (cross) term, we prove by
induction on the power $m$ that
\begin{equation}\label{eq:cross-term}
  \norm{(a+b)^m - a^m - b^m}
    \le (\norm{a}+\norm{b})^m - \norm{a}^m - \norm{b}^m
    \qquad (m \ge 1),
\end{equation}
using the non-commutative recurrence
$(a+b)^{m+1} - a^{m+1} - b^{m+1} = (a+b)\bigl((a+b)^m - a^m - b^m\bigr) + ab^m + ba^m$.
Summing over the exponential series, the cross term is also bounded by
$(e^{\norm{a}}-1)(e^{\norm{b}}-1)$, so by the triangle inequality:
\begin{equation}\label{eq:step-bound}
  \norm{e^a\,e^b - e^{a+b}}
    \le 2(e^{\norm{a}}-1)(e^{\norm{b}}-1)
    \le 2\norm{a}\norm{b}\,e^{\norm{a}+\norm{b}},
\end{equation}
where the last step uses $(e^s-1)(e^t-1) \le st\,e^{s+t}$ for $s,t \ge 0$.

\begin{remark}
  The factorization \eqref{eq:factorization} was discovered during
  formalization: the standard second-order expansion approach required
  tracking non-commutative multinomial cross-terms, which was intractable
  in Lean. The algebraic factorization reduces the problem to two
  independent bounds, each handled by series comparison.
\end{remark}

\subsection{Assembly and convergence}
\label{sec:assembly}

Setting $P = e^{A/n}\,e^{B/n}$ and $Q = e^{(A+B)/n}$:
\begin{enumerate}
  \item $Q^n = e^{A+B}$ by the identity $(e^{a/n})^n = e^a$
    (\leanref{ExpDivPow.lean}{31}{exp\_div\_pow}).
  \item $\norm{P-Q} \le 2\norm{A}\norm{B}/n^2 \cdot e^{(\norm{A}+\norm{B})/n}$
    by the step error with $a = A/n$, $b = B/n$.
  \item $\max\{\norm{P},\norm{Q}\} \le e^{(\norm{A}+\norm{B})/n}$ by
    \eqref{eq:norm-exp}.
  \item Telescoping gives
    $\norm{P^n - Q^n} \le n \cdot \norm{P-Q} \cdot M^{n-1}$.
  \item The factors $e^{(\norm{A}+\norm{B})/n}$ combine:
    $e^{s/n} \cdot (e^{s/n})^{n-1} = (e^{s/n})^n = e^s$, yielding
    \[
      \norm{(e^{A/n}\,e^{B/n})^n - e^{A+B}}
        \le \frac{2\norm{A}\norm{B}\,e^{\norm{A}+\norm{B}}}{n}.
    \]
\end{enumerate}
The main theorem $\lim_{n\to\infty} (e^{A/n}\,e^{B/n})^n = e^{A+B}$ follows
by a standard $\varepsilon$-$N$ argument using the Archimedean property.

\subsection{Duhamel approach to commutator scaling}
\label{sec:commutator}

The commutator-scaling bound \eqref{eq:comm-first} replaces the product
$\norm{A}\norm{B}$ with $\norm{\comm{B}{A}}$. Our proof uses a
variation-of-parameters (Duhamel) technique based on the fundamental theorem
of calculus, avoiding ODE uniqueness theory entirely.

\paragraph{Integral error representation.}
Define the ``conjugated Trotter product''
\begin{equation}\label{eq:w-def}
  w(\tau) = e^{-\tau(A+B)}\,e^{\tau B}\,e^{\tau A}.
\end{equation}
By the product rule for Fr\'{e}chet derivatives in a normed algebra,
\begin{equation}\label{eq:w-deriv}
  w'(\tau) = e^{-\tau(A+B)}\,\bigl(e^{\tau B}\,A - A\,e^{\tau B}\bigr)\,e^{\tau A}
  = e^{-\tau(A+B)}\,\comm{e^{\tau B}}{A}\,e^{\tau A}.
\end{equation}
The key cancellation is that the ``free'' terms involving $A+B$, $B$, and $A$
sum to zero: $-(A+B) + B + A = 0$. The FTC-2 gives
\begin{equation}\label{eq:integral-error}
  e^{tB}\,e^{tA} - e^{t(A+B)}
    = \int_0^t e^{(t-\tau)(A+B)}\,\comm{e^{\tau B}}{A}\,e^{\tau A}\,d\tau,
\end{equation}
after left-multiplying by $e^{t(A+B)}$ (which commutes with itself).

\paragraph{Extracting the commutator.}
The integrand contains $\comm{e^{\tau B}}{A} = e^{\tau B}A - Ae^{\tau B}$,
which involves the \emph{exponential} of $B$. To extract the bare commutator
$\comm{B}{A}$, we apply FTC-2 a second time to the conjugation map
$s \mapsto e^{sB}\,A\,e^{-sB}$:
\begin{equation}\label{eq:conj-ftc}
  e^{\tau B}\,A\,e^{-\tau B} - A
    = \int_0^{\tau} e^{sB}\,\comm{B}{A}\,e^{-sB}\,ds.
\end{equation}
This gives $\norm{e^{\tau B}\,A\,e^{-\tau B} - A}
  \le \norm{\comm{B}{A}}\,\abs{\tau}\,e^{2\abs{\tau}\norm{B}}$
by bounding the integrand and using
$\norm{\int_0^\tau f} \le C\abs{\tau}$.
Right-multiplying by $e^{\tau B}$ yields
$\norm{\comm{e^{\tau B}}{A}} \le \norm{\comm{B}{A}}\,\abs{\tau}\,e^{3\abs{\tau}\norm{B}}$.

\paragraph{Final bound.}
Substituting into \eqref{eq:integral-error} and bounding the integral:
\[
  \norm{e^{tB}\,e^{tA} - e^{t(A+B)}}
    \le \int_0^t e^{(t-\tau)(\norm{A}+\norm{B})} \cdot
      \norm{\comm{B}{A}}\,\tau\,e^{3\tau\norm{B}} \cdot e^{\tau\norm{A}}\,d\tau.
\]
Simplifying the exponential factors and evaluating $\int_0^t \tau\,d\tau = t^2/2$
gives the claimed bound \eqref{eq:comm-first}.

\begin{remark}
  The Duhamel approach requires only the product rule for Fr\'{e}chet
  derivatives, the FTC-2 for Bochner integrals, and basic norm estimates---all
  available in \Mathlib. The alternative approach via Gronwall's inequality
  would require ODE existence/uniqueness theory, adding approximately 40 lines
  of additional Lean infrastructure.
\end{remark}

\subsection{Second-order Strang commutator scaling}
\label{sec:strang}

For the Strang splitting $S_2(t) = e^{tA/2}\,e^{tB}\,e^{tA/2}$, we prove
\begin{equation}\label{eq:strang-bound}
  \norm{S_2(t) - e^{tH}}
    \le \frac{\norm{\comm{B}{\comm{B}{A}}}}{12}\,t^3
      + \frac{\norm{\comm{A}{\comm{A}{B}}}}{24}\,t^3,
\end{equation}
where $H = A + B$ and we additionally require that $A$ and $B$ are
\emph{skew-adjoint} ($a^* = -a$), so that $\norm{e^{ta}} = 1$.

The proof applies the same Duhamel strategy to
$w(\tau) = e^{-\tau H}\,S_2(\tau)$, but now requires:
\begin{enumerate}
  \item A \emph{4-factor product rule} for $S_2(\tau) = e^{\tau A/2}\,e^{\tau B}\,e^{\tau A/2}$,
    yielding $w'(\tau) = e^{-\tau H}\,\mathcal{T}_2(\tau)\,S_2(\tau)$, where $\mathcal{T}_2$
    involves the commutators $\comm{A/2}{e^{\tau B}}$ and $\comm{e^{\tau B}}{A/2}$.

  \item \emph{First-order cancellation}: the leading commutator terms cancel
    exactly: $\comm{A/2}{B} + \comm{B}{A/2} = 0$.

  \item A \emph{double FTC} on $s \mapsto e^{sB}\,A\,e^{-sB}$, expanding to
    second order to extract $\comm{B}{\comm{B}{A}}$.

  \item The \emph{``subtract-constant-at-$\tau$'' trick}: to bound
    $\int_0^\tau F(s)\,ds - \tau\,F(\tau)$, rewrite as
    $\int_0^\tau (F(s) - F(\tau))\,ds$ and bound $\norm{F(s) - F(\tau)}
    \le C_A(\tau - s)$.
    This avoids Fubini's theorem or integration by parts.
\end{enumerate}

\paragraph{The 4-factor product rule.}
Write $A' = A/2$ and $H = A + B$. The conjugated Strang product
\[
  w(\tau) = e^{-\tau H}\,e^{\tau A'}\,e^{\tau B}\,e^{\tau A'}
\]
has four exponential factors, each contributing a derivative term via the
Leibniz rule. Using $d/d\tau[e^{\tau X}] = X\,e^{\tau X}$ and collecting:
\begin{equation}\label{eq:4factor}
  w'(\tau) = e^{-\tau H}\,\mathcal{T}_2(\tau)\,S_2(\tau),
\end{equation}
where the \emph{Strang residual} is
\[
  \mathcal{T}_2(\tau)
    = \bigl(e^{\tau A'}\,B\,e^{-\tau A'} - B\bigr)
      + e^{\tau A'}\bigl(e^{\tau B}\,A'\,e^{-\tau B} - A'\bigr)e^{-\tau A'}.
\]
The key algebraic step is verifying that the ``free'' derivative terms
$-H + A' + B + A'$ sum to zero (since $A' + A' = A$ and $H = A + B$).
In Lean, the sum of the four derivative contributions is factored as
$-(e^{-\tau H}) \cdot (H - A' - A' - B) \cdot S_2(\tau)$, where
$H - A' - A' - B = 0$ is discharged by \texttt{abel}. The free-ring
factoring is handled by \texttt{noncomm\_ring}.

\paragraph{Coefficient conversion.}
Since $A' = A/2$, the double commutators in the final bound involve
half-scaled operators:
$\comm{B}{\comm{B}{A'}} = \frac{1}{2}\comm{B}{\comm{B}{A}}$ and
$\comm{A'}{\comm{A'}{B}} = \frac{1}{4}\comm{A}{\comm{A}{B}}$.
The scalar factors are extracted through $\comm{cX}{Y} = c\,\comm{X}{Y}$
(valid for central scalars) using \texttt{Algebra.smul\_mul\_assoc} and
\texttt{Algebra.mul\_smul\_comm}, then \texttt{norm\_smul}.
Combined with the integration constants ($1/6$ from the cubic integral
and $1/2$ from the double FTC), this produces the final coefficients
$1/12$ and $1/24$ in \eqref{eq:strang-bound}.
